\section{Tendências Tecnológicas Emergentes}

A arquitetura dos gêmeos digitais na próxima década será redefinida pela convergência de fluxos de dados contínuos e inteligência artificial \textit{multi-omics}. 

\subsection{A Transição para Gêmeos Digitais Dinâmicos (Tempo Real)}

A primeira fronteira tecnológica disruptiva é a ruptura com os modelos estáticos. Atualmente, o gêmeo digital é uma "fotografia tridimensional" extraída no momento do escaneamento. O futuro pertence às representações dinâmicas alimentadas por telemetria contínua (\textit{Digital Twin 2.0}). A incorporação de biossensores intraorais — encapsulados em alinhadores ortodônticos, implantes \textit{smart} ou placas de bruxismo — permitirá o monitoramento de tensores de força (\textit{strain gauges}), oscilações de pH, biomarcadores salivares e perfis de microbiota em tempo real \cite{khoshfekr2025digital}. 

Essa conectividade (IoT Odontológica) transforma o cuidado episódico em monitoramento ativo. Se o sensor em um implante acusar micro-movimentações que precedem a perda óssea, o Gêmeo Digital recalcula imediatamente o vetor de estresse oclusal e emite um alerta preditivo, permitindo a intervenção cirúrgica profilática. Esse nível de antecipação inaugura a era da ``odontologia preventiva guiada por dados'' \cite{shen2024effectiveness}.

\subsection{Fusão Multimodal Avançada e Fenotipagem Profunda}

O segundo vetor é a \textbf{Fenotipagem Profunda} (\textit{Deep Phenotyping}). A biomecânica isolada é insuficiente para explicar falhas como a peri-implantite crônica. O gêmeo digital do futuro integrará, em uma malha holística, a geometria craniofacial (CBCT), o transcriptoma do ligamento periodontal e o microbioma gengival \cite{techsoc2025towards}. Redes neurais de \textit{cross-attention} conseguirão correlacionar uma variação geométrica da crista óssea com a expressão de marcadores inflamatórios específicos, prescrevendo não apenas o posicionamento do implante, mas a modulação farmacológica necessária para aquele paciente específico.

\subsection{A Era Prescritiva: Do 'O Que Acontecerá?' para 'O Que Devo Fazer?'}

Modelos preditivos informam ao dentista que uma coroa de zircônia irá fraturar sob determinada carga. Modelos prescritivos vão além: eles executam automaticamente milhares de simulações de Monte Carlo no plano de fundo (background) e recomendam, de forma autônoma, uma alteração no ângulo da cúspide palatina que mitigará o risco de fratura para menos de 0.1\%, sugerindo inclusive a resina ideal para a cimentação \cite{menon2026digital}. A transição do Gêmeo Diagnóstico para o Gêmeo Prescritivo consolida o sistema como um \textit{Co-Piloto} clínico.
